\documentclass[conference]{IEEEtran}
\IEEEoverridecommandlockouts
\usepackage{cite}
\usepackage{amsmath,amssymb,amsfonts}
\usepackage{array}
\usepackage{booktabs}
\usepackage{graphicx}
\usepackage{textcomp}
\usepackage{xcolor}
\usepackage{url}
\newcolumntype{L}[1]{>{\raggedright\arraybackslash}p{#1}}
\emergencystretch=2em
\def\BibTeX{{\rm B\kern-.05em{\sc i\kern-.025em b}\kern-.08em
    T\kern-.1667em\lower.7ex\hbox{E}\kern-.125emX}}

\begin{document}

\title{Beyond Film Thickness: A Design-Only Research Plan for Correlative Geometric--Molecular--Force Metrology of Nanoscale Film Stability}

\author{\IEEEauthorblockN{Anonymous Research Plan Author}
\IEEEauthorblockA{\textit{Design-only research-planning artifact}\\
No institutional affiliation asserted\\
Non-biological interface scope only}}

\maketitle

\begin{abstract}
Microscopic interfaces govern drainage, wetting, reaction, and interfacial transport, yet a nanometre-scale film-thickness trace does not identify the molecular or mechanical causes of a transition. This paper integrates a manual literature Survey, a hypothesis-generation artifact, and a safety-gated ExperimentDesign artifact into a design-only research plan for non-biological liquid--solid thin-film interfaces. The central proposition is conditional: within an expert-approved, geometry-matched comparison boundary, a temporally coherent relation among film geometry, an interface-specific molecular-state representation, and a local confined-interaction representation may discriminate chemistry-mediated disjoining-pressure control of drainage or stability from a thickness-only explanation.

The contribution is an evidence architecture, not an empirical finding. The plan specifies a geometry-only baseline; complementary geometric, molecular, and local-force representations; a cross-channel registration model; a structured alternative set; and four outcome branches that include an explicit no-information state. It identifies decisions that cannot be automated from the supplied evidence: the model interface, chemical-safety and facility review, calibration traceability, comparability boundary, estimand, and replication rationale. The manuscript supplies no material identity, chemical formulation, apparatus configuration, numerical operating parameter, physical procedure, or observed result. It is intended to prevent a later correlation from being retrospectively treated as a causal mechanism.
\end{abstract}

\begin{IEEEkeywords}
interfacial metrology, nanoscale liquid films, disjoining pressure, sum-frequency generation, confined forces, uncertainty registration, causal discrimination, research design
\end{IEEEkeywords}

\section{Introduction}

An interface is not adequately described by its location alone. At gas--liquid, liquid--solid, and solid--solid boundaries, a small change in separation can coincide with molecular reorientation, adsorption, solvation restructuring, capillary deformation, confined reaction, heat transfer, or mass transfer. A measurement that resolves only one of these descriptions cannot, in general, determine which one explains a transition. The practical result is a mismatch between the observable that is easiest to obtain and the mechanism that is most important to understand.

Optical interference and ellipsometric methods are central to microscopic-interface research because they can provide non-contact geometric information about thin films. The ability to resolve thickness, curvature, or lateral morphology is indispensable. It nevertheless does not uniquely identify molecular orientation, local chemical state, a confined interaction landscape, or the route by which a film drains, arrests, persists, or ruptures. The Survey underlying this paper therefore begins from a non-identifiability claim: a thickness record is necessary for many questions, but insufficient for a causal explanation.

The literature provides complementary component capabilities. Nonlinear optical spectroscopy can constrain interfacial molecular and orientational states~\cite{backus2020,tang2020}. Quantitative interpretation of a sum-frequency vibrational representation requires explicit spectral, orientational, polarization, and model assumptions~\cite{wang2005}; this limitation is scientifically useful because it prevents a molecular signature from being treated as a direct causal readout. Surface-sensitive vibrational methods can also address selected buried interfaces~\cite{chen2010}. At the mechanical scale, scanning-probe and surface-force methods can constrain nanoscale morphology, separation, capillary objects, and confined interactions~\cite{lohse2015,dziadkowiec2019}. Imaging ellipsometry provides a non-contact geometric anchor for nanometre-scale liquid films~\cite{shoji2024}. Thermal metrology can indicate whether an interface state matters for transport, but a transport response alone is often spatially averaged and mechanistically indirect~\cite{cahill2003}.

These capabilities motivate a correlative architecture. They do not establish that a combined geometric--molecular--force relation has already caused a stability transition in a particular system. This paper preserves the distinction between component-method evidence and a novel causal hypothesis. It offers a research plan rather than an empirical report, applies only to future expert-approved non-biological liquid--solid model interfaces, and authorizes no physical work.

\section{Literature Synthesis and the Measurement Gap}

\subsection{What an interface measurement must distinguish}

The phrase ``measuring an interface'' conceals several different scientific tasks. Geometry concerns thickness, curvature, roughness, and morphology. Molecular state concerns orientation, adsorption, hydration, coordination, and local chemical state. Mechanics concerns confined force, capillary interaction, effective potential, and a possible disjoining-pressure contribution. Transport concerns evaporation, heat flow, reaction, or mass flux. Each task has different observables, spatial footprints, time response, and perturbation modes.

The interface literature explains why a single modality does not close all tasks. Broad reviews identify coupled structural, chemical, and dynamic open questions rather than one universal descriptor~\cite{bjorneholm2016}. Nonlinear optical spectroscopy is useful when interfacial molecular arrangement is the object of interest~\cite{backus2020}, but the signal-to-structure relation is conditional on experimental geometry and analysis assumptions~\cite{wang2005}. Surface-force evidence addresses a different question: whether a confined interaction changes with separation and surface state~\cite{dziadkowiec2019}. Optical geometry can be less perturbative and more continuous in time, yet it remains dependent on refractive-index and optical-model assumptions~\cite{shoji2024}.

The target of this proposal is not a claim that every interface requires every modality. It is a framework for deciding when a causal explanation of a thin-film transition requires evidence beyond geometry. The central question is whether complementary channels are mutually constraining. If they are not, a geometry-only or measurement-artifact explanation must remain viable.

\subsection{Component capabilities and retained limits}

Table~\ref{tab:modalities} distinguishes what a channel can constrain from what it cannot independently establish. The table is neither a protocol nor an instrument-selection guide. It records the evidentiary role of each family and the limitation that remains live during interpretation.

\begin{table*}[!t]
\caption{Complementary measurement roles and retained interpretive limits.}
\label{tab:modalities}
\centering
\footnotesize
\begin{tabular}{L{0.20\textwidth}L{0.23\textwidth}L{0.25\textwidth}L{0.22\textwidth}}
\toprule
Evidence family & Primary representation in this plan & Contribution to causal discrimination & Limit that remains active \\
\midrule
Optical interference and ellipsometry & Time-indexed thickness, curvature, or morphology & Supplies the geometric baseline and tests whether candidate accounts are geometry matched & Refractive-index, optical-model, and lateral-heterogeneity assumptions can distort the represented state \\
Surface-specific nonlinear vibrational spectroscopy & Interface-specific molecular or orientational state & Constrains whether a molecular-state transition accompanies a geometric transition & Spectral inversion, phase, polarization, and orientational assumptions prevent direct unqualified causal attribution \\
Scanning-probe and surface-force approaches & Local topography, separation, or confined-force proxy & Constrains whether an interaction landscape changes in a mechanism-consistent manner & Probe perturbation, roughness, confinement, and interaction-model error can create an apparent force relation \\
Chemical-state channels & Selected adsorbate or chemical-state constraint & Can strengthen an interface-state alternative when compatible with the comparison boundary & Information depth, enhancement substrate, scattering, and observation perturbation are method specific \\
Thermal or transport-sensitive metrology & Transport-consequence representation & Tests whether a proposed interface state matters for a transport outcome & A transport signal does not uniquely identify a molecular cause \\
\bottomrule
\end{tabular}
\end{table*}

The primary architecture uses the first three roles---geometry, interface-specific molecular state, and local confined interaction---because each addresses a different part of a candidate mechanism. A transport channel is deliberately optional. The Survey contains useful thermal-transport context, but it does not supply a sufficiently direct cross-validated cluster for one reactive or evaporative application regime. Restricting this plan to drainage or stability avoids overextending the evidence base.

\subsection{Precise research gaps}

The literature synthesis identifies three linked gaps. \textbf{G1: thickness-only non-identifiability.} A thickness or morphology representation cannot independently identify molecular state, local force, or cause. \textbf{G2: temporal and spatial registration.} Complementary channels can observe related states at different positions, times, or degrees of perturbation; a loose cross-channel correlation is not causal evidence. \textbf{G3: perturbation and model audit.} Surface history, roughness, contamination, probe effects, interaction models, and observation-induced chemistry can produce an apparently coherent story. These factors must be explicit rival explanatory paths rather than caveats added after a preferred conclusion.

The contribution of the present proposal is to treat these gaps jointly. It does not claim a new measurement technology. It specifies a decision architecture in which the evidence favoring a mechanism and the evidence capable of invalidating that favor are both first-class outputs.

\section{Research Objective, Scope, and Central Proposition}

\subsection{Research objective}

The objective is to formulate a falsifiable, auditable plan for deciding whether a future non-biological liquid--solid thin-film comparison can discriminate chemistry-mediated local-interaction control of drainage or stability from a thickness-only account. A bounded model interface is to be selected later by qualified experts. The paper does not identify a liquid, solid, surface treatment, chemical additive, sample source, apparatus, or operating condition.

The question is causal rather than merely correlative: when a transition is represented, what additional evidence is necessary before part of that transition can be attributed to a molecular-state-dependent confined interaction? The plan gives a conditional evidence requirement, not a guarantee that the required evidence can always be acquired.

\subsection{Central proposition}

The central proposition is: \emph{within a bounded non-biological liquid--solid thin-film system, a temporally coherent relation among film geometry, an interface-specific vibrational or chemical-state representation, and a confined-force proxy can distinguish chemistry-mediated disjoining-pressure control of drainage or stability from a thickness-only explanation.}

For planning purposes, define the represented state at relative event time $t$ as
\begin{equation}
\mathbf{X}(t)=\bigl[G(t),M(t),F(t),T(t),A(t)\bigr],
\label{eq:state-vector}
\end{equation}
where $G(t)$ is calibrated geometry, $M(t)$ is the interface-specific molecular-state representation, $F(t)$ is the local confined-interaction representation, $T(t)$ is the declared drainage or stability-transition outcome, and $A(t)$ is the audit record. The audit record includes calibration assumptions, surface history, perturbation assessment, interaction-model assumptions, and temporal-registration uncertainty.

Equation~\eqref{eq:state-vector} is descriptive. It does not assert that entries are measured on a common scale, are simultaneous, or are sufficient to close a causal mechanism. Its purpose is to prevent an incomplete evidence record from being represented as complete merely because several signals have been collected.

The competing explanatory accounts are
\begin{equation}
\mathcal{M}_{0}:T\leftarrow(G,A),\qquad
\mathcal{M}_{1}:T\leftarrow(G,M,F,A),
\label{eq:model-comparison}
\end{equation}
where $\mathcal{M}_{0}$ is a geometry-plus-audit baseline and $\mathcal{M}_{1}$ is the joint account. The notation defines information supplied to each account; it does not prescribe a statistical model, numerical threshold, or physical procedure. Any future estimator must be selected only after the model system, uncertainty structure, and estimand are approved.

\subsection{Scope and falsifier}

The proposal is not a general theory of all disjoining-pressure phenomena and is not a design for evaporative cooling, spray drying, combustion, or materials synthesis. Those applications motivate the importance of microscopic-interface metrology but require focused evidence and safety review beyond the supplied corpus. The paper makes no universal claim across liquids, solids, roughness regimes, surface chemistries, or interface geometries.

The proposition is falsified or narrowed if a geometry-plus-audit account predicts the transition equally well; if molecular and force representations do not add coherent discriminating information; if apparent timing fails under registration uncertainty; or if optical calibration, surface history, capillarity, roughness, probe effects, interaction-model error, or observation-induced chemistry explains the relation. These are not secondary caveats. They define the hypothesis itself.

\section{Conceptual Model of Cross-Channel Comparability}

\subsection{Registration is part of the scientific claim}

Cross-channel comparison is often treated as a data-integration task after physical interpretation has been chosen. In this plan it is part of the causal claim. If a molecular channel observes a state at a different effective time or location from the geometric and force channels, then an apparent sequence can be a property of observation rather than of the interface.

For channel $k$, a generic registered representation is
\begin{equation}
\widetilde{X}_{k}(t)=X_{k}\bigl(t+\delta_{k}(t)\bigr)+\varepsilon_{k}(t),
\label{eq:registration}
\end{equation}
where $\widetilde{X}_{k}$ is the retained representation, $X_{k}$ is the latent channel-specific state, $\delta_{k}(t)$ is timing or registration offset, and $\varepsilon_{k}(t)$ represents measurement and model uncertainty. Equation~\eqref{eq:registration} is not an instruction to estimate these terms from a particular dataset. It identifies what must be bounded before a statement such as ``molecular change precedes the transition'' can be meaningful.

The comparability boundary has four dimensions. \emph{Geometry comparability} asks whether channels refer to adequately matched separation, curvature, morphology, and lateral state. \emph{Temporal comparability} asks whether effective observation windows have a documented relation. \emph{Surface-history comparability} asks whether the interface has an auditable preparation and exposure history. \emph{Perturbation comparability} asks whether observation changes the phenomenon it observes. The boundary is successful only if it is restrictive enough to invalidate non-comparable records and transparent enough to be audited.

\subsection{Molecular and force states as constraints}

The molecular-state channel is a constraint on which accounts remain plausible, not a mechanism label. A future plan would have to state what aspect of the molecular representation is relevant, which interpretation assumptions support that mapping, and which alternatives remain unresolved. If this cannot be done, the molecular channel may remain descriptive but cannot carry the causal weight assigned to it in $\mathcal{M}_{1}$.

The same rule applies to local force. A confined-force observation can be affected by probe geometry, roughness, solvent confinement, approach history, and the interaction model used to map an observed response to a physical quantity. The plan treats $F(t)$ as a proxy representation; it cannot be relabeled as a direct disjoining-pressure term unless the mapping is justified in a future approved system. The mechanism requires consistency, not identity, between $M(t)$ and $F(t)$. A molecular change unrelated to local interaction may be irrelevant to stability, while a force change without molecular support may arise from geometry or probe artifacts.

\section{Design Architecture}

\subsection{Logical comparison accounts}

The design uses four logical accounts. They are not material recipes or experimental samples; they determine what an eventual evidence record must distinguish.

\begin{table}[!t]
\caption{Logical accounts for mechanism discrimination.}
\label{tab:accounts}
\centering
\footnotesize
\begin{tabular}{@{}L{0.13\columnwidth}L{0.36\columnwidth}L{0.39\columnwidth}@{}}
\toprule
Account & Included representations & Decision role \\
\midrule
C0 & $G,T,A$ & Geometry-plus-audit baseline: tests whether the transition is adequately explained without molecular or force additions \\
C1 & $G,M,T,A$ & Molecular-added account: asks whether molecular state adds information beyond C0 \\
C2 & $G,F,T,A$ & Interaction-added account: asks whether local interaction adds information beyond C0 \\
C3 & $G,M,F,T,A$ & Joint account: asks whether molecular and force representations jointly improve discrimination over C0--C2 \\
\bottomrule
\end{tabular}
\end{table}

The decisive comparison is not C3 versus a favorable trace. It is whether C3 provides a coherent, audit-resilient explanation beyond C0 and the two single-added-channel accounts. For a future approved analysis, the incremental discriminating value can be written conceptually as
\begin{equation}
\Delta_{\mathrm{disc}}=\mathcal{Q}(\mathcal{M}_{1};T\mid A)-\mathcal{Q}(\mathcal{M}_{0};T\mid A),
\label{eq:discrimination}
\end{equation}
where $\mathcal{Q}$ denotes an expert-approved measure of explanatory adequacy conditional on the audit record. The plan leaves $\mathcal{Q}$ unspecified. Choosing an estimand, likelihood, score, or threshold before the model interface and uncertainty structure are known would create false precision. Equation~\eqref{eq:discrimination} records the comparative burden: added channels must make a difference beyond geometry while remaining robust to the audit.

\subsection{Required audit records}

Table~\ref{tab:audit} identifies potential routes by which an apparent mechanism can fail. A future study may add records, but it may not drop a listed record and still classify an inference as strong causal support.

\begin{table*}[!t]
\caption{Audit requirements and invalidation logic.}
\label{tab:audit}
\centering
\footnotesize
\begin{tabular}{L{0.22\textwidth}L{0.30\textwidth}L{0.34\textwidth}}
\toprule
Audit domain & Question before interpretation & Consequence if unresolved \\
\midrule
Geometric calibration & Does the geometry representation retain model assumptions and uncertainty across the comparison boundary? & Geometry may be misrepresented; no added-channel claim is interpretable \\
Surface history and roughness & Are history, contamination risk, and roughness sufficiently documented to compare states? & An apparent molecular or force shift can remain a surface-history explanation \\
Channel perturbation & Can the observation channel or probe alter chemistry, morphology, or interaction? & The represented transition can remain measurement induced \\
Interaction-model adequacy & Is the relation between local response and force proxy explicitly bounded? & The force channel cannot support a disjoining-pressure interpretation \\
Temporal and spatial registration & Is effective location and time basis documented with uncertainty? & Apparent precedence or correlation can be an alignment artifact \\
Missing or invalid channel evidence & Is every evidence role available inside the comparability boundary? & Select the uninformative-or-invalid branch; do not impute mechanism support \\
\bottomrule
\end{tabular}
\end{table*}

\subsection{Analysis plan without false operational precision}

No sample size, repeat count, allocation scheme, instrument configuration, material identity, or statistical threshold is set here. These omissions are intentional. A repetition rationale requires a defined experimental unit, uncertainty model, state variability, and estimand. A calibration plan requires a specific model system and compatible observation channels. An instrument configuration requires facility and safety review. Filling these fields now would make the plan appear more executable while reducing its scientific honesty.

The proposal is nevertheless methodologically specific. A future approved design must predeclare a transition outcome and geometry-only baseline; state the interpretation of molecular and force proxies; retain calibration and registration uncertainty; distinguish independent repetitions from repeated representations of one state; record exclusions and unavailable channels; and apply the conditional outcome branches. These requirements constrain future work without becoming a laboratory protocol.

\section{Uncertainty, Alternatives, and Falsification}

\subsection{Alternative-explanation matrix}

The target mechanism is one member of an alternative set. Table~\ref{tab:alternatives} makes that set explicit. A later report should not merely say that an alternative was considered; it must state which representation or audit bears on it and whether that evidence can discriminate it.

\begin{table*}[!t]
\caption{Alternative explanations that remain live until discriminated.}
\label{tab:alternatives}
\centering
\footnotesize
\begin{tabular}{L{0.21\textwidth}L{0.30\textwidth}L{0.35\textwidth}}
\toprule
Alternative explanation & Why it can mimic the target relation & Evidence requirement for a bounded conclusion \\
\midrule
Geometry or capillarity alone & A transition can follow curvature, separation, or capillary forcing without a molecular-state mechanism & An explicit geometry-plus-audit baseline compared with added-channel accounts \\
Optical-model or refractive-index error & Represented thickness or timing can change because the optical mapping changes & Stated calibration model, uncertainty, and sensitivity to justified alternatives \\
Surface history, contamination, or roughness & These factors can change molecular and force signals without the proposed state transition & Traceable surface-history record and declared comparability boundary \\
Probe perturbation or interaction-model error & A local observation can create, alter, or mis-map the interaction being interpreted & Perturbation assessment and bounded proxy interpretation \\
Observation-induced chemistry or confinement & The measurement environment can change adsorption, reaction, or dynamics & Explicit assessment of channel-induced state changes where appropriate \\
Temporal or spatial mismatch & Non-coincident records can create apparent order or coupling & Registration model and uncertainty adequate to test timing claims \\
\bottomrule
\end{tabular}
\end{table*}

\subsection{Falsifier and conditional outcome branches}

The main falsifier is deliberately strong: the chemistry-mediated account is not supported if the geometry-plus-audit account predicts the declared transition equally well, or if no temporally coherent molecular-plus-force relation remains after the alternatives are assessed. The proposal is also not supported if a force proxy cannot be interpreted inside its model limits or if the molecular channel remains non-specific in the chosen comparison boundary.

Falsification is not a failure of the workflow. A geometry-only result would clarify that added channels supplied no discriminating information in the documented regime. An invalid result would clarify that the comparison boundary or audit record was inadequate. Both outcomes are more useful than an uninterpretable affirmative claim.

\begin{table*}[!t]
\caption{Conditional outcome branches and permitted claims.}
\label{tab:branches}
\centering
\footnotesize
\begin{tabular}{L{0.18\textwidth}L{0.29\textwidth}L{0.25\textwidth}L{0.18\textwidth}}
\toprule
Branch & Future trigger & Permitted interpretation & Required next action \\
\midrule
Supports mechanism & Joint evidence adds coherent explanatory information beyond geometry and the audit does not favor an alternative & Bounded evidence would be consistent with chemistry-mediated local interaction in the documented model boundary only & Independent review and replication before any broader claim \\
Partial or heterogeneous & The relation is informative only in a subset of approved states or is boundary sensitive & The relation may be conditional or incompletely identifiable & Refine state taxonomy and inspect heterogeneity sources \\
Null or contradictory & Geometry explains equally well, or added channels conflict & The proposed mechanism is not supported in the documented boundary & Report the null or contradiction; revise the Idea separately if warranted \\
Uninformative or invalid & A critical calibration, comparability, perturbation, or timing audit is absent or non-discriminating & No causal conclusion is permitted & Stop interpretation and redesign the evidence architecture \\
\bottomrule
\end{tabular}
\end{table*}

\section{Data Governance and Reproducibility}

The planned work is reproducible first at the level of reasoning. The paper preserves the proposition, baseline account, required representations, alternatives, outcome branches, and open decisions. A later reviewer can therefore determine whether a future report has changed the hypothesis after observing a favorable signal.

If a physical comparison is later approved, the reproducibility record must retain raw and derived channel representations; calibration metadata; time-registration decisions; surface-history metadata; uncertainty assumptions; analysis versions; exclusions; unavailable observations; and the rationale for branch selection. Each record is necessary to reconstruct why an alternative was retained, bounded, or favored.

The source lineage is qualified. The upstream Survey was manually authored using cross-validated bibliographic records, but it has no automatic Survey-to-Idea manifest. The Idea and ExperimentDesign artifacts therefore retain the status \texttt{survey\_\allowbreak unbound\_\allowbreak requires\_\allowbreak human\_\allowbreak confirmation}. The status does not invalidate the synthesis; it prevents manual provenance from being represented as an automatically verified handoff.

\section{Responsible Research and Human Review}

The ExperimentDesign artifact has execution mode \texttt{DESIGN\_ONLY} and validation status \texttt{BLOCKED\_\allowbreak BY\_\allowbreak RISK\_\allowbreak REVIEW}. This paper neither authorizes nor describes physical work. It includes no human participants, animals, plants, clinical populations, environmental samples, microbial systems, cell lines, or restricted biological materials. It supplies no chemical formulation, material identity, apparatus configuration, operating parameter, quantity, or stepwise procedure.

The following decisions remain \emph{needs human input}: selection of a scientifically relevant and safe non-biological model interface; chemical-safety and facility confirmation; compatibility of geometric, molecular, and local-force representations; calibration and time-registration traceability; the meaning of a transition endpoint; the estimand and uncertainty model; and an independent-repetition rationale. Reviewers may decide that a literature-only evidence map, a formal causal model, or no further action is preferable to a physical route. The proposal remains valid under any of these choices because it is a decision architecture rather than an execution plan.

\section{Expected Contribution and Limitations}

The expected contribution is methodological specificity. The plan converts the broad instruction ``measure the interface more completely'' into finite scientific obligations. It identifies a baseline that must be beaten, channels that must be interpreted rather than merely acquired, and audits that can invalidate an apparent mechanism. It also separates component-method evidence from the novel combined hypothesis. This is important in interdisciplinary interface science, where individually credible measurements can be combined into a persuasive narrative without actually referring to the same physical state.

The central limitation is equally clear: the paper has no observed results. It cannot select the best model interface, establish cross-channel compatibility, quantify uncertainty, or prove transferability. The component literature spans different interfaces and methodological contexts. A future result in one approved system would remain bounded and could not be generalized automatically to another liquid, solid, or application. The manual provenance status also requires confirmation before an automated pipeline claim is made.

These limitations are productive. They define next decisions and preserve falsifiability. The plan rejects two weak outcomes: a broad but untestable assertion that chemistry matters, and a detailed but unauthorized procedure whose parameters appear precise without scientific or safety justification.

\section{Conclusion}

This paper presents a design-only plan for correlative geometric--molecular--force metrology of nanoscale liquid-film stability. It does not assert that chemistry-mediated disjoining-pressure control has been established. Instead, it defines evidence conditions under which a future, non-biological, expert-approved comparison could distinguish that account from geometry-only explanations, identify conditionality, report a null or contradiction, or correctly conclude that evidence is non-interpretable.

The architecture is deliberately conservative: geometry anchors the state; molecular and local-force representations add potentially discriminating information; and calibration, surface history, perturbation, interaction-model, and registration audits determine whether that information can carry a causal interpretation. Making these dependencies explicit before execution provides a stronger foundation for human review than a thickness record or an unregistered cross-channel correlation alone.

\section*{Acknowledgment}

This document is a design-only planning artifact. It reports no experimental data and declares no funding, institutional affiliation, or execution authorization.

\begin{thebibliography}{00}

\bibitem{bjorneholm2016} O. Bj\"orneholm \emph{et al.}, ``Water at interfaces,'' \emph{Chemical Reviews}, 2016, doi: 10.1021/acs.chemrev.6b00045.
\bibitem{backus2020} E. H. G. Backus, J. Schaefer, and M. Bonn, ``Probing the mineral-water interface with nonlinear optical spectroscopy,'' \emph{Angewandte Chemie International Edition}, 2020, doi: 10.1002/anie.202003085.
\bibitem{tang2020} F. Tang, T. Ohto, S. Sun, \emph{et al.}, ``Molecular structure and modeling of water-air and ice-air interfaces monitored by sum-frequency generation,'' \emph{Chemical Reviews}, 2020, doi: 10.1021/acs.chemrev.9b00512.
\bibitem{wang2005} H. Wang, W. Gan, R. Lu, Y. Rao, and B. Wu, ``Quantitative spectral and orientational analysis in surface sum frequency generation vibrational spectroscopy (SFG-VS),'' \emph{International Reviews in Physical Chemistry}, 2005, doi: 10.1080/01442350500225894.
\bibitem{chen2010} Z. Chen, ``Investigating buried polymer interfaces using sum frequency generation vibrational spectroscopy,'' \emph{Progress in Polymer Science}, 2010, doi: 10.1016/j.progpolymsci.2010.07.003.
\bibitem{lohse2015} D. Lohse and X. Zhang, ``Surface nanobubbles and nanodroplets,'' \emph{Reviews of Modern Physics}, 2015, doi: 10.1103/RevModPhys.87.981.
\bibitem{dziadkowiec2019} J. Dziadkowiec, B. Zareeipolgardani, D. K. Dysthe, and A. R\o yne, ``Nucleation in confinement generates long-range repulsion between rough calcite surfaces,'' \emph{Scientific Reports}, 2019, doi: 10.1038/s41598-019-45163-6.
\bibitem{shoji2024} E. Shoji, A. Hoshino, T. Biwa, \emph{et al.}, ``Superspreading wetting of nanofluid droplet laden with highly dispersed nanoparticles,'' \emph{Langmuir}, 2024, doi: 10.1021/acs.langmuir.4c03347.
\bibitem{cahill2003} D. G. Cahill, W. K. Ford, K. E. Goodson, \emph{et al.}, ``Nanoscale thermal transport,'' \emph{Journal of Applied Physics}, 2003, doi: 10.1063/1.1524305.

\end{thebibliography}

\end{document}
